\section{Conclusion: The Dialectics of Formalization}

The trajectory from extensional structures through conceptualizations to ontologies enacts a fundamental epistemological progression. We begin with the mathematical bedrock—sets and relations—devoid of modal or intensional content. We enrich this structure by introducing possible worlds, obtaining conceptual relations as functions that remain invariant across modal variation. We commit our formal languages to these intensional structures, determining intended models. Finally, we attempt to axiomatically approximate the space of intended models through ontologies.

At each stage, we confront limitations:
\begin{itemize}
	\item Extensional structures cannot distinguish essential from accidental relations
	\item Conceptualizations, while intensionally rich, cannot be directly expressed in formal languages
	\item Ontological commitments determine intended models but do not uniquely identify them
	\item Ontologies approximate intended models but cannot fully capture conceptualizations
\end{itemize}

Yet these limitations are not mere deficiencies but rather constitutive features of formal representation itself. As von Bertalanffy recognized, ``the whole is more than the sum of its parts'' precisely because constitutive organization transcends extensional enumeration. The intensional turn—the introduction of conceptual relations defined over possible worlds—is our response to this recognition.

The ontology, as explicit specification, becomes the site where language meets reality, where formal rigor confronts the irreducible richness of conceptual content. It is approximation not by failure but by necessity—the necessary incompleteness of any finite axiomatization before the infinite modal space of conceptualization.

In the interstices between formal precision and conceptual adequacy, we find the living practice of ontological engineering: an ongoing negotiation between what can be said and what must be shown, between the extensional clarity of logic and the intensional depth of understanding.